\subsection{Driver Models}
\subsubsection{Intelligent Driver Model (IDM)}

The Intelligent Driver Model (IDM) is a deterministic microscopic car-following model describing the longitudinal dynamics of vehicle $i$ as a function of its own velocity, the spacing to its predecessor, and their relative velocity \cite{Treiber.2000}. The acceleration of vehicle $i$ is defined as

\begin{equation}
	\dot{v}_i =
	a_{max} \left[
	\underbrace{1 - \left( \frac{v_i}{v_d} \right)^{\delta}}_{\text{free-road acceleration}}
	-
	\underbrace{\left( \frac{s^*(v_i,\Delta v_i)}{s_i} \right)^2}_{\text{interaction (braking) term}}
	\right].
\end{equation}
Here $v_i$ is the ego velocity, $v_d$ the desired free-flow velocity, $a_{max}$ the maximum acceleration, $b$ the desired deceleration, and $\delta$ the acceleration exponent \cite{Treiber.2000}. The free-road term drives the vehicle toward $v_d$, while the interaction term enforces deceleration when the actual spacing $s_i$ falls below the desired spacing $s^*$.


The spacing is defined as
\begin{equation}
	s_i = x_{i-1} - x_i - L,
\end{equation}

where $x_i$ and $x_{i-1}$ denote the longitudinal positions of vehicle $i$ and its leader, respectively, and $L$ is the vehicle length. The relative velocity is defined as $\Delta v_i = v_i - v_{i-1}$.

The desired dynamic spacing is given by

\begin{equation}
	s^*(v_i,\Delta v_i)
	=
	\underbrace{s_0}_{\text{jam distance}}
	+
	\underbrace{v_i T_H}_{\text{time-headway term}}
	+
	\underbrace{\frac{v_i \Delta v_i}{2 \sqrt{a_{max} b}}}_{\text{anticipation term}}.
\end{equation}
The terms represent the jam distance $s_0$, the time-headway policy $v_i T_H$, and an anticipative braking correction \cite{Treiber.2000}.


\subsubsection{Stochastic Intelligent Driver Model (SIDM)} \label{SIDM}

To account for perception errors, behavioral variability, and unmodeled disturbances, the IDM can be extended by introducing a stochastic acceleration component \cite{Treiber.2018}. The stochastic IDM (SIDM) is formulated as

\begin{equation}
	\dot{v}_i =
	\underbrace{f(s_i, v_i, v_{i-1})}_{\text{deterministic IDM}}
	+
	\underbrace{\xi_i(t)}_{\text{stochastic acceleration noise}},
\end{equation}

where $f(\cdot)$ denotes the deterministic IDM acceleration function and $\xi_i(t)$ is a stochastic process representing white acceleration noise.

The stochastic term satisfies
\begin{equation}
	\langle \xi_i(t) \rangle = 0,
\end{equation}
\begin{equation}
	\langle \xi_i(t)\xi_j(t') \rangle
	=
	Q \, \delta_{ij} \, \delta(t - t'),
\end{equation}

where $Q$ denotes the noise intensity, $\delta_{ij}$ is the Kronecker delta ensuring independence between different vehicles, and $\delta(t-t')$ is the Dirac delta distribution representing temporal white noise \cite{Treiber.2018}.

 Close to the string stability threshold, noise can generate spatiotemporal correlations and traffic oscillations resembling stop-and-go waves. Consequently, the SIDM provides a unified framework for analyzing deterministic instabilities and noise-induced traffic dynamics \cite{Treiber.2018}.